% Source: https://artofproblemsolving.com/community/c3h1634418_how_to_do_this_without_trig
% Stable identity: aops-topic-1634418-q01
Name the triangle $ABC$ with $\angle A=15^\circ$, $\angle B=45^\circ$ and $\angle C=120^\circ$; the side opposite $B$ is $AC=20$. Drop the perpendicular from $A$ to line $BC$; since $\angle C$ is obtuse, its foot $H$ lies on the extension of $BC$ beyond $C$.

\begin{asy}
usepackage("amsmath");
size(6.5cm);
defaultpen(fontsize(10pt));
real s=sqrt(3);
pair H=(0,0),A=(0,10*s),C=(-10,0),B=(-10*s,0);
draw(A--B--C--cycle,linewidth(0.9));
draw(C--H,dashed);draw(A--H,dashed);
draw((-1.2,0)--(-1.2,1.2)--(0,1.2));
draw(arc(B,2.5,0,45),blue);label("$45^\circ$",B+3.7*dir(15),blue);
draw(arc(C,1.6,60,180),blue);label("$120^\circ$",C+3.0*dir(105),blue);
draw(arc(C,2.2,0,60),red);label("$60^\circ$",C+3.3*dir(28),red);
label("$A$",A,N);label("$B$",B,SW);label("$C$",C,S);label("$H$",H,SE);
label("$20$",(A+C)/2,E);
label("\mbox{}",(min(currentpicture,user=true)+max(currentpicture,user=true))/2,black); // keeps the text after the figure black when LaTeX typesets the labels inline
\end{asy}

\textbf{Triangle $ACH$.} The exterior angle at $C$ is $\angle ACH=180^\circ-120^\circ=60^\circ$, so $ACH$ is a $30$-$60$-$90$ triangle with hypotenuse $AC=20$: $CH=10$ and $AH=10\sqrt3$.

\textbf{Triangle $ABH$.} It has a right angle at $H$ and $\angle B=45^\circ$, so it is isosceles: $BH=AH=10\sqrt3$. Hence $BC=BH-CH=10\sqrt3-10$.

\textbf{Area.} Taking $BC$ as the base, the height is $AH$:
\[
[ABC]=\dfrac{1}{2}\cdot BC\cdot AH=\dfrac{1}{2}\left(10\sqrt3-10\right)\cdot10\sqrt3=\dfrac{300-100\sqrt3}{2}=150-50\sqrt3 .
\]
So $m=150$, $n=50$, $q=3$, and $m+n+q=\boxed{203}$.
